\documentclass[11pt]{article}
\usepackage{fontspec}
\setmainfont{Times New Roman}
\setsansfont{Arial}
\setmonofont{Consolas}
\usepackage{amsmath,amssymb}
\usepackage{geometry}
\usepackage{microtype}
\geometry{a4paper,margin=3cm}
\setlength{\parindent}{1.25em}
\setlength{\parskip}{0pt}
\emergencystretch=30em
\allowdisplaybreaks
\raggedbottom
\providecommand{\mM}{\mathfrak M}
\providecommand{\mG}{\mathfrak G}
\providecommand{\mU}{\mathfrak U}
\providecommand{\mB}{\mathfrak B}
\providecommand{\ma}{\mathfrak a}
\providecommand{\mb}{\mathfrak b}
\providecommand{\me}{\mathfrak e}
\providecommand{\mx}{\mathfrak x}

\begin{document}

% Унит: Папер 17 Сецтион 08. Дирецт Интерславиц Латин транслатион фром те Герман финал-аудитед соурце слице.

\setcounter{footnote}{16}
\subsection*{\S~8. Адитивны разклад пополно редуцибилных груп до простых груп. Теорем о изоморфији.}

Ныне хочемо разсмотрити другы случај, в ктором сице не владаје једнозначность, але владаје изоморфија разних разкладов. Ту јде о обобченье случаја, размотреного од Лоевы.

\textbf{Дефиниција.} \textit{Модул называ се простым модулом, ако не има других делительев кроме самого себе и јединичного модула, обхвачајучего все полиномы. Група остатков простого модула называ се проста група.}

Всака проста група јест а фортиори иредуцибилна. Обставајут такоже бесконечне просте групы (сравни дефиницију \S~11 и примјер \S~12, 2).

Наша дана група остатков \(\mG\) неха буде ныне представима на два способы како сума (всак раз конечно многих) простых груп \(\mU_1,\ldots,\mU_k\) односно \(\mB_1,\ldots,\mB_l\). Групу остатков, ктора допушча најмене једно тако представјенье како суму простых груп, и јеј одповедајучи модул \(\mM\) называјемо, следујучи Лоевыјевој означбе, пополно редуцибилными.

Ныне разматрјамо продукты \(\ma_\chi\mb_\lambda\), кде \(\chi=1,\ldots,k\), \(\lambda=1,\ldots,l\), и хочемо показати, же или \(\ma_\chi\mb_\lambda=0\), или \(\mx\ma_\chi\mb_\lambda\) пробегаје целу групу \(\mB_\lambda\), когда \(\mx\) пробегаје все класы остатков од \(\mG\). В самој ствари, модул належечи к \(\mB_\lambda\)
\[
 (\mM,\mb_1+\cdots+\mb_{\lambda-1}+\mb_{\lambda+1}+\cdots+\mb_l)=\mM_{\mb_\lambda},
\]
кторы настаје из \(\mM\), когда додајемо все полиномы класа остатков \(\mb_1+
\cdots+\mb_{\lambda-1}+\mb_{\lambda+1}+\cdots+\mb_l\), има делитель
\[
 (\mM,\mb_1+\cdots+\mb_{\lambda-1}+\mb_{\lambda+1}+\cdots+\mb_l,
 \ma_\chi\mb_\lambda).
\]
Але понеже модул \(\mM_{\mb_\lambda}\) по предположеньу јест просты модул, тој делитель јест или равны \(\mM_{\mb_\lambda}\), или јест јединичны модул. В првом случају \(\ma_\chi\mb_\lambda\) належи к \(\mM_{\mb_\lambda}\); зато обставаје релација
\[
 \ma_\chi\mb_\lambda=\mx(\mb_1+
\cdots+\mb_{\lambda-1}+\mb_{\lambda+1}+\cdots+\mb_l),
\]
из чего следује \(\ma_\chi\mb_\lambda=0\). В другом случају обставајут два класа \(\mx_0\) и \(\mx_1\), тако же
\[
 \mx_0\ma_\chi\mb_\lambda
 +\mx_1(\mb_1+\cdots+\mb_{\lambda-1}+\mb_{\lambda+1}+\cdots+\mb_l)
 =\me=\mb_1+\cdots+\mb_{\lambda-1}+\mb_{\lambda+1}+\cdots+\mb_l+\mb_\lambda,
\]
и зато, заради једнозначности, \(\mx_0\ma_\chi\mb_\lambda=\mb_\lambda\). Зато \(F\mx_0\ma_\chi\mb_\lambda\), и теды такоже \(\mx\ma_\chi\mb_\lambda\), пробегаје целу групу \(\mB_\lambda\), когда \(F\) односно \(\mx\) пробегаје все полиномы односно все класы остатков од \(\mG\). Тым тврдженье јест доказано.

При \(\ma_\chi\mb_\lambda\ne0\) јест далье \(\mU_\chi\) изоморфна к \(\mB_\lambda\). Бо целокупност всех класов \(\mx\ma_\chi\ne0\), за кторе такоже \(\mx\ma_\chi\mb_\lambda\ne0\), јест по \S~7 изоморфна к системе класов \(\mx\ma_\chi\mb_\lambda\); а ты изчрпавајут подгрупу \(\mB_\lambda\). Треба теды само показати, же именоване класы \(\mx\ma_\chi\) изчрпавајут групу \(\mU_\chi\). За то разматрјамо все класы остатков \(\mx\ma_\chi\), за кторе \(\mx\ma_\chi\mb_\lambda=0\). Тута система има својство, же сума двох јеј класов и продукт с либовольным полиномом знову належи к системе. Ако зато додамо все те класы остатков к модулу \(\mM_{\ma_\chi}\), настаје знову модул, кторы јест делитель од \(\mM_{\ma_\chi}\), теды или \(\mM_{\ma_\chi}\), или јединичны модул. В последньем случају бы вшак сама \(\ma_\chi\) была в ньем садржеча, теды \(\ma_\chi\mb_\lambda=0\) против предположеньа. Зато не обставајут класы \(\mx\ma_\chi\ne0\), за кторе \(\mx\ma_\chi\mb_\lambda=0\), и изоморфија од \(\mU_\chi\) к \(\mB_\lambda\) јест доказана.

Неха ныне за некоје определьено \(\chi\), на примјер,
\[
 \ma_\chi\mb_{\lambda_1}\ne0,
 \ldots,\ma_\chi\mb_{\lambda_i}\ne0\qquad(i>1).
\]
Тогда \(\mU_\chi\) јест изоморфна к групам \(\mB_{\lambda_1},\ldots,\mB_{\lambda_i}\), зато ты групы сут меджу собоју изоморфне.

Ныне хочемо показати, же тогда и најмене две из груп \(\mU\) сут изоморфне. Бо ако бы в продуктах \(\mb_\lambda\ma_\chi\) к всакому \(\mb_\lambda\) належала само једна \(\ma_\chi\), тако же \(\mb_\lambda\ma_\chi\ne0\), теды ако бы \(\mb_\lambda=\mb_\lambda\ma_\chi\), тогда бы было \(\mB_\lambda=\mU_\chi\), понеже \(\mU_\chi\) како проста група не може садрживати властну подгрупу. Тогда бы было \(k=l\) и при правилном нумерованьу \(\ma_\chi=\mb_\chi\); але тогда бы такоже схема \(\ma_\chi\mb_\lambda\) была диагонална схема, что не јест случај. Тым доказан јест следујучи теорем:

\textbf{Теорем V.} \textit{Ако пополно редуцибилна група \(\mG\) јест на два способы представльена како сума простых груп:
\[
 \mG=\mU_1+\cdots+\mU_k=\mB_1+\cdots+\mB_l,
\]
тогда всака група \(\mU\) јест изоморфна најмене једнеј групе \(\mB\), и обратно. Ако всака \(\mU\) јест изоморфна точно једнеј \(\mB\), тогда разклады сут идентичне. В противном случају најмене две групы \(\mU\) сут меджу собоју изоморфне, и такоже најмене две групы \(\mB\) сут меджу собоју изоморфне.}

\end{document}

